\section{Source interface: supported zero and original norms}
The Zeta workbench's SplitZero construction \cite{Derived,Tau,Toda} is
used here, not replaced by a second scalar theory. For a ring $k$, its
coefficient semiring is $G(k)=\{\tau\}\sqcup k^\bullet$, with
$e=0_k^\bullet\ne\tau$. Reconstruction over a support semilattice $L$
forms a disjoint union of module fibres. Its internal quotient is
\[
 (\lambda,x)\longmapsto(\lambda,[x]);
 \qquad (\lambda,0)\ne\tau\quad(\lambda\ne\bot).
\]
The difference between two representatives is an actual relation, not a
reason to remove the fibre label. In the source's tau-base diagram the
cochain is $[V\oplus V\xrightarrow{r_+-r_-}\mathscr B]$ in degrees $0,1$.
The analytic theta maps, their domains and $\mathscr B/\Theta V$ are not
identified with a finite ES module in this note. We instantiate the
source's \emph{algebraic} reconstruction, internal homology, kernel and
original-Gram quotient interfaces on specified finite arithmetic maps.
No purity, RH, or analytic-source comparison is imported.

The source's least-norm construction is likewise retained. If $G>0$ is
the original Gram, $B$ the independent relation columns and $C$ a chosen
representative matrix, then
\[
 R=C-B(B^*GB)^{-1}B^*GC,\quad
 R^*GR=C^*GC-C^*GB(B^*GB)^{-1}B^*GC.                 \tag{1}
\]
The determinant of the source Gram in columns $[C,B]$ is
$\det(B^*GB)\det(R^*GR)$. This is the existing Schur-complement calculation
\cite[\S1]{Toda}, not a new determinant theorem. Our contribution below
is a complete arithmetic instance, its counting weights and original
selector inverse. The finite ES norm is specified directly; it is not
silently declared equal to the theta-function norm.

\section{The complete finite cochain and its quotient}
Let $X,Y,Z$ be finite sets with maps $\ell:X\to Z$, $r:Y\to Z$.
For a commutative ring $k$, use the free modules with their labelled bases.
Let $X_z=\ell^{-1}(z)$ and $Y_z=r^{-1}(z)$, and choose one reference
$x_z$ or $y_z$ in each nonempty fibre. The relations are
\[
 B_X=\bigoplus_{z}\bigoplus_{x\in X_z\setminus\{x_z\}}k(e_x-e_{x_z}),
 \qquad B_Y=\bigoplus_{z}\bigoplus_{y\in Y_z\setminus\{y_z\}}k(e_y-e_{y_z}).
\]
Consider the actual three-term cochain
\[
 \mathcal C^{-1}=B_X\oplus B_Y
 \xrightarrow{\partial}\mathcal C^0=k^X\oplus k^Y
 \xrightarrow{d}\mathcal C^1=k^Z,
 \quad d(a,b)_z=\sum_{X_z}a_x-\sum_{Y_z}b_y.       \tag{2}
\]
The incoming map is inclusion; $d\partial=0$ term by term.

\begin{theorem}[Complete integral cohomology and inverse]\label{thm:cohom}
Put $S=\ell(X)$, $T=r(Y)$ and $I=S\cap T$. Over every commutative ring,
\[
 H^{-1}(\mathcal C)=0,\quad H^0(\mathcal C)\simeq k^I,
 \quad H^1(\mathcal C)\simeq k^{Z\setminus(S\cup T)}.            \tag{3}
\]
The degree-zero isomorphism sends a cycle to
$c_z=\sum_{X_z}a_x=\sum_{Y_z}b_y$. Its inverse sends $c_z$ to
$c_z(e_{x_z},e_{y_z})$. The complete fibre over $(c_z)$ consists of this
representative plus $\partial\theta$, with unique relation coefficients
$\theta_x=a_x$ and $\theta_y=b_y$ away from the reference indices.
\end{theorem}
\begin{proof}
The displayed star differences are independent and generate exactly the
kernel of each fibre sum: a vector of sum zero equals
$\sum_{x\ne x_z}a_x(e_x-e_{x_z})$. In a cycle the two sums agree. If
only one fibre is nonempty, this common sum is zero; otherwise it is any
$c_z$. Subtracting the stated reference representative leaves precisely
the displayed relations. This proves (3), including both inverse
compositions, without division in $k$. Finally the image of $d$ contains
exactly all coordinate basis vectors indexed by $S\cup T$.
\end{proof}

Taking only $\ker d$ without the original incoming relations would count
internal cancellations as witnesses. The relation primitive in the proof
specifies exactly what is removed. To apply the source's support theory,
use the four masks $L=\mathcal P(\{\mathrm{left},\mathrm{right}\})$,
coordinate inclusions at present legs, and the common codomain for every
nonempty mask. Set the empty-mask complex to zero. The relation submodules
are transition-stable. Reconstruct (2) and its homology internally as in
\cite[\S\S1--4]{Derived}; the differential-square equation is the
supported-zero map at the retained mask, not the constant absent map.
If the joint $H^0$ is zero, that is still its \emph{present zero fibre}.

\subsection{The complete canonical weight}
Give $k^X\oplus k^Y$ its original real or complex diagonal Gram, with
positive entries $a_x,b_y$. Define
\[
 M_z=\sum_{X_z}a_x^{-1},\qquad N_z=\sum_{Y_z}b_y^{-1}.
\]
For $z\in I$, its least-norm representative is
\[
 r(c)_x=\frac{c_z}{a_xM_z},\qquad
 r(c)_y=\frac{c_z}{b_yN_z},                            \tag{4}
\]
and its quotient norm is
\[
 \boxed{\|c\|_{H^0}^2=\sum_{z\in I}
          (M_z^{-1}+N_z^{-1})|c_z|^2.}                \tag{5}
\]
To prove this, the weighted inner product of (4) with each star relation
is zero. Its sums are $c_z$, and its squared norm is (5). Subtracting it
from any other cycle supplies the explicit relation primitive and the
Pythagorean identity. Empty fibres are excluded from (4), not assigned a
fictitious inverse weight. Least-norm representatives may be rational or
complex; integral witness extraction uses labelled occurrences, not the
claim that (4) is an integral source vector.

For the original counting Gram $a_x=b_y=1$, put $m_z=|X_z|$, $n_z=|Y_z|$.
Then the quotient Gram is $\operatorname{diag}(1/m_z+1/n_z)$.
In the integral cycle basis given by the reference representative and
star differences, a common fibre has source determinant $m_z+n_z$,
relation determinant $m_zn_z$, and quotient determinant
$(m_z+n_z)/(m_zn_z)$. A one-sided fibre contributes the same determinant
to source and relations. These formulas follow either from (1) or from
$\det(I+\mathbf1\mathbf1^*)=\text{matrix size}+1$ on each star block.
Thus all source and relation factors are visible.

The canonical projection of the original all-ones vector onto the
harmonic degree-zero subspace has components
$2n_z/(m_z+n_z)$ on $X_z$ and $2m_z/(m_z+n_z)$ on $Y_z$, with squared norm
\[
 \boxed{4\sum_{z\in I}\frac{m_zn_z}{m_z+n_z}.}       \tag{6}
\]
This is positive exactly when $I$ is nonempty. It is an exact source-norm
criterion, not a bound proving $I$ nonempty for arbitrary arithmetic data.

\subsection{Coarsening and killed classes}
For $\pi:Z\to Z'$, retain the same $X,Y$ but replace their maps by
$\pi\ell,\pi r$. The identity on degree zero and fibre-sum map on degree
one define a chain map; the incoming map sends each old star difference
to its explicit expression in the coarser star basis. The induced map is
\[
 (c_z)_{z\in I}\longmapsto
 \left(\sum_{\substack{z\in I\\\pi(z)=w}}c_z\right)_{w\in I'},
 \quad I'=\pi(S)\cap\pi(T).                         \tag{7}
\]
Its image is $k^{\pi(I)}$, its cokernel $k^{I'\setminus\pi(I)}$, and its
kernel has the star-difference basis in each fibre of $I\to\pi(I)$.
For any killed class its chain representative has coarser sums zero, and
the preceding nonreference-coordinate formula is its boundary primitive.
This is an actual instance of the source's transported-class kernel,
not the assumption that comparisons are injective \cite[\S3]{Derived}.
For example two matched residues merging to one kill the nonzero class
$(1,-1)$; unmatched fine residues can also create new coarse matches.

\section{The ES specialization: two channels, all actual exponents}
In this arithmetic application, the coefficient field is $\mathbb Q$.
Fix a prime $p\equiv1\pmod4$, $p/4<a<p/2$, $R=4a-p$, and
$a=\prod_iq_i^{e_i}$. Put $B_e=\prod_i[-e_i,e_i]_{\mathbb Z}$ and
$v(\beta)=\prod_iq_i^{\beta_i}\pmod R$. All powers are units.
Retain the two \emph{canonical} channel targets
\[
 t_M=-1,\qquad t_E=-p^{-1}\pmod R.                  \tag{8}
\]
The additional inverted-exterior target of the split-energy test is a
labelled alternative computation of E, not a third canonical channel.

Split $e=f+b$ with $f,b\ge0$. In each separately labelled channel use
\[
 X=B_f,\quad Y=B_b,\quad Z=(\mathbb Z/R\mathbb Z)^\times,
 \quad \ell(\xi)=v(\xi),\quad r_c(\eta)=t_cv(\eta)^{-1}.       \tag{9}
\]
The full object is the direct sum of these channel complexes, retaining
$(p,a,R,c,f,b)$ as its type. Theorem~\ref{thm:cohom} proves
\[
 H^0(\mathcal C_c)\ne0
 \iff \exists(\xi,\eta):v(\xi+\eta)=t_c.             \tag{10}
\]
Every $\beta\in B_e$ has the complete addition fibre
\[
 \max(-f_i,\beta_i-b_i)\le\xi_i\le\min(f_i,\beta_i+b_i),
 \qquad \eta=\beta-\xi,                             \tag{11}
\]
so (10) is equivalent to an original channel hit. The original divisor is
\[
 u=\prod_iq_i^{e_i+\beta_i},\quad d_0=\gcd(a,u),\quad
 (h,r,s)=(d_0^2/u,u/d_0,a/d_0).                       \tag{12}
\]
On M use $\lambda=(r+s)/R$ and denominators
$(a,phs\lambda,phr\lambda)$; on E use $\kappa=(pr+s)/R$ and
$(a,hs\kappa,phr\kappa)$. The inverse is
$u=pa^2/(RY-pa)$ on M and $u=a^2/(RY-pa)$ on E.
Prime valuations prove all normalizations; clearing the reciprocal sum
proves the original ES identity. These coordinates and first-half
completeness are inherited \cite{ET,ESBase,Split}. The ordered pair
$(\xi,\eta)$, channel and full exponents are not inferred from a residue
alone; they are selected from the recorded fibres (9),(11).

\section{The exact original arithmetic weight and its chain primitive}
Let $\mathscr P_c$ be the matched pairs in (9) and let
$\mathscr W_c=\{\beta\in B_e:v(\beta)=t_c\}$. Addition gives a surjection
\[
 q:\mathscr P_c\to\mathscr W_c,
 \qquad N(\beta)=|q^{-1}(\beta)|
 =\prod_i\bigl(\min(f_i,\beta_i+b_i)-\max(-f_i,\beta_i-b_i)+1\bigr).\tag{13}
\]
There are \emph{two} observations of $\mathscr P_c$: its matched residue
and its original exponent vector. Neither determines the other in general.
The residue cohomology rank, the matching-pair count and the original
witness count must not be identified.

For any finite surjection $q:P\to W$, let $Q:k^P\to k^W$ sum over its
fibres. Its kernel has the complete basis $e_x-e_{x_0(w)}$. Thus
\[
 [\ker Q\hookrightarrow k^P]\simeq k^W[0]          \tag{14}
\]
in degrees $-1,0$, by the integral reference section and explicit fibre relations. Over
$\mathbb Q$ the canonical section for the counting Gram is
\[
 J(c)_x=c_{q(x)}/N(q(x)),\quad QJ=1.                 \tag{15}
\]
For every amplitude vector $A$, retain both $QA$ and
$\theta_x=A_x-J(QA)_x$ at each nonreference point. This is a linear
bijection, with inverse $J(QA)+\sum\theta_x(e_x-e_{x_0})$.
In particular a nonzero cancellation is an actual boundary with its
primitive, not external absence.

The inherited counting norm gives
\[
 \|A\|^2=\sum_w\frac{|QA(w)|^2}{N(w)}+\|A-JQA\|^2.             \tag{16}
\]
If instead one is transporting the \emph{original uniform norm} on
$k^W$, its permutation-invariant diagonal metric on a refined fibre is
necessarily $N(w)I$: minimization then gives quotient norm $|c_w|^2$.
This transported metric is displayed as a changed metric; it is not called
the native product-counting norm. Indeed a constant diagonal $h_wI$ on a
fibre has quotient norm $h_w|c_w|^2/N(w)$, proving uniqueness.

Over $\mathbb Q$ or $\mathbb R$, arbitrary positive original ES masses $\omega_{c,\beta}$ are
reconstructed without duplication by assigning to each lifted pair the
mass $\omega_{c,\beta}/N(\beta)$. In particular the full original
channel count and every monomial are recovered exactly:
\[
 \boxed{\sum_{\beta\in\mathscr W_c}U^{e+\beta}
 =\sum_{(\xi,\eta)\in\mathscr P_c}
          \frac{U^{e+\xi+\eta}}{N(\xi+\eta)}.}      \tag{17}
\]
The zero coefficient is absence of that monomial; a signed intermediate
sum cancelling to zero retains its source support mask separately.

\subsection{All failed candidates retain their original defect}
For every split pair, matched or not, put
\[
 x=v(\xi),\quad w=v(\eta),\quad\rho=x-t_cw^{-1}\pmod R,
 \qquad \beta=\xi+\eta.
\]
The original gates $g_M=u+a$, $g_E=4u+1$ satisfy
\[
 g_c\equiv c_c\rho\pmod R,
 \qquad c_M=aw,\quad c_E=pw,                        \tag{18}
\]
where both multipliers are units. Hence
\[
 D_c(\beta)=\frac R{\gcd(R,g_c)}=\frac R{\gcd(R,\rho)}.         \tag{19}
\]
For the original affine monodromy
\[
 (X,Y,Z)\mapsto(6g+Y,-6g-X-Y,Z-2g+X),\quad
 (-Y,X-6g,Z+3g+Y),\quad(X,Y+X,Z-g),
\]
scaling by $c_c$ is an equivariant bijection from coefficient $\rho$ to
coefficient $g_c$, inverse scaling by $c_c^{-1}$. No global sphere theorem
is required. Equations (13),(18) prove the \emph{complete defect-polynomial}
identity
\[
 \sum_{\beta\in B_e}U^{e+\beta}Z^{D_c(\beta)-1}
 =\sum_{\xi\in B_f,\eta\in B_b}
       \frac{U^{e+\xi+\eta}Z^{D_c(\xi+\eta)-1}}{N(\xi+\eta)}. \tag{20}
\]
Evaluating the polynomial at $Z=0$ means coefficient extraction at exponent
zero and gives (17). This is not a limiting positivity assumption.

\subsection{Iterated quotients must carry their inherited metric}
For finite surjections $P\xrightarrow q W\xrightarrow r Z$, let
$m_w=|q^{-1}(w)|$ and $n_z=\sum_{r(w)=z}m_w$.
The quotient metric on $k^W$ is $\operatorname{diag}(1/m_w)$.
The canonical lift of $c_z$ to $w$ is $m_wc_z/n_z$; lifting again gives
$c_z/n_z$ at every original occurrence. It agrees exactly with the direct
canonical lift for $rq$. Every kernel coordinate is recovered by the
successive star primitives. This proves all finite towers by induction.

Resetting the intermediate metric to the identity changes that map.
For three budgets $[-1,1]$, total exponent zero has seven ordered triples.
The direct coefficients are all $1/7$ and squared norm $1/7$.
Resetting after adding the first two gives four coefficients $1/6$ and
three $1/9$, with squared norm $4/27$. Thus the exact excess is
\[
 \boxed{\frac4{27}-\frac17=\frac1{189}>0.}           \tag{21}
\]
The discarded difference is a nonzero zero-sum boundary. This is a concrete
reason to retain the source/relation metric through each morphism.

\section{Concrete tests and the counterfactual object}
At $p=2521$, $a=636=2^2\cdot3\cdot53$, $R=23$, choose
$f=(1,1,1)$, $b=(1,0,0)$. The complete matched pair table is
\[
\begin{array}{c|r|r|r|r|r}
c&g&\xi&\eta&\beta&N(\beta)\\\hline
M&22&(-1,1,1)&(0,0,0)&(-1,1,1)&2\\
M&11&(0,-1,-1)&(1,0,0)&(1,-1,-1)&2\\
M&21&(0,1,1)&(-1,0,0)&(-1,1,1)&2\\
M&22&(1,-1,-1)&(0,0,0)&(1,-1,-1)&2\\
E&13&(-1,1,0)&(-1,0,0)&(-2,1,0)&1
\end{array}
\]
There are five matched pairs, four channel/residue cohomology coordinates,
and three original canonical channel states. The harmonic Gram in the
ordered shared coordinates
\[ (M,11),\ (M,21),\ (M,22),\ (E,13) \]
is $\operatorname{diag}(2,2,3/2,2)$, determinant $12$. The mass identity is
$4\cdot(1/2)+1=3$, not five. The two middle states have
$(u,h,r,s,\lambda)=(8,2,2,159,7)$ and $(50562,2,159,2,7)$; the exterior
state has $(477,53,3,4,329)$. Equation (12) returns all denominators and
(17) retains the two middle orientations as distinct canonical labels.

At the genuinely empty \emph{single shell} $p=241,a=64,R=15$, split
$6=3+3$. The two channel complexes together have dimensions
\[
 \dim C^{-1}=12,\quad \dim C^0=28,\quad\dim C^1=16,
 \qquad H^{-1}=H^0=H^1=0.
\]
The raw two-leg kernel has dimension twelve; all twelve coordinates are
actual relations. The counting norm of the all-ones source vector is
$28$, not zero. Its supported acyclic fibre is not $\tau$. Other shells
at 241 succeed, so this is not an ES counterexample.

Likewise at $p=37,a=18,R=35$, the exterior local divisor sets are
$\{1,6,36,81\}$ modulo five and $\{12,54\}$ modulo seven. They are both
nonempty but disjoint. Equation (7) permits new cohomology after
coarsening; it does not give a simultaneous original exponent vector.
One must retain the joint residue or the same original leaf in every local
factor. These original regression sets are inherited \cite{ESBase}.

For the complete first-half atlas at a fixed prime, choose any split at
each shell and sum the two canonical channel complexes with their full
support labels. Its degree-zero cohomology is nonzero exactly when ES
holds at that prime. A hypothetical false prime has zero ordinary
$H^0$ in every such labelled summand, not absent source objects or zero
source weights. On the refined matched-pair complexes (14), $H^0$ counts
actual canonical states; on (2) it counts matched channel/residues.
The distinction persists under all split changes.

These identities do not produce a universal nonzero class. They supply
an exact source-norm expression (6), a complete original state-count
expression (17), and an all-defect expression (20) on which a positivity
or cross-shell estimate can be attempted. A generic source-norm bound from
the Zeta application is not presumed valid for this finite arithmetic
family merely because the quotient formulas have the same algebraic form.

\section{Verification scope}
The standard-library verifier constructs the actual differentials, relation
columns and canonical representatives. It checks their ranks, both inverse
compositions, relation primitives, and determinant ratios for every pair
of maps from sets of size at most two into a three-element set, in both
counting and specified nonconstant diagonal metrics. It also checks unequal
common fibres through size six, every small split budget, and the complete
three-stage tower. The displayed arithmetic examples include every original
candidate and every split pair, successful or failed. Further bounded
first-half scans retain their exact domain in the JSON. They prove no
unbounded occupancy statement.

Source-level general statements and finite checks have different scopes.
The Zeta analytic source is read through the pinned workbenches cited below;
the full 482-page reader and the analytic/formal corpus were not audited or
replayed. No Lean theorem is claimed for this tranche. Existing mathematical
source files are not rewritten, and no historical-priority claim is made.
